\section{Exponential finite-record complexity and reconstruction}
\label{sec:exponential-reconstruction}
The finite graph used in Lemma~\ref{lem:finite-record-variation}
has both its number of variables and its number of polynomial tests
linear in the collision length. Retaining the binomial form of the
component bound, rather than its fixed-dimension asymptotic estimate,
gives an exponential variation bound. This removes the additional
logarithm from the regularized tower budget.

\subsection{The finite binomial bound}
\begin{lemma}[Linear-size sign descriptions]
\label{lem:linear-sign-complexity}
Fix $a_0,a_1,d\ge1$. Suppose a family of lifted graphs of length $L$
uses at most $a_0(L+1)$ real variables and $a_1(L+1)$ polynomial
sign tests of degrees at most $d$, and has at most $a_0^{L+1}$
discrete alternatives. The number of connected components of any
Boolean combination of those sign conditions, summed over the
alternatives, is at most $C e^{CL}$. The same bound holds after
coordinate projection and after fixing a bounded number of coordinates
or adjoining a bounded number of fixed-degree tests. The constants
depend on $a_0,a_1,d$, not on the real coefficients.
\end{lemma}
\begin{proof}
Use the finite bound in \cite[Section 3.2, (3.3)]{BPR}, with zeroth
Betti number and ambient variety the full space. If $s$ is the number
of tests and $k\ge1$ the number of variables, it bounds the sum of
the numbers of connected components of all realizable sign conditions
by
\begin{equation}\label{eq:finite-binomial-exponential}
 d(2d-1)^{k-1}\sum_{j=0}^{k}\binom{s}{j}4^j
 \le d(2d-1)^{k-1}5^s.
\end{equation}
Here $\binom{s}{j}=0$ for $j>s$. The right side is exponential in
$k+s$, so it is $Ce^{CL}$ in the stated range. A Boolean combination
is a union of realizable sign conditions. Each of its connected
components contains a connected component of one of those conditions;
thus its number is no larger than their total. Summing over the
exponentially many discrete alternatives preserves an exponential bound.

The image of a connected set under a coordinate projection is connected.
A union of $m$ connected sets has at most $m$ connected components.
This proves the projection assertion without using quantitative
elimination of quantifiers. Fixed coordinates and additional tests
only change the linear bounds by constants. A bounding-ball inequality
can likewise be adjoined. In zero dimension the assertion is immediate.
We use the finite inequality in \eqref{eq:finite-binomial-exponential},
not an $s\to\infty$ expansion with constants depending on $k$.
\end{proof}

\begin{theorem}[Exponential variation of physical finite records]
\label{thm:exponential-finite-record}
There is $C_{\mathrm g}\ge1$, uniform in $R\in I$, such that all the
conclusions of Lemma~\ref{lem:finite-record-variation} hold with
\begin{equation}\label{eq:exponential-record-budget}
 B_L=C_{\mathrm g}e^{C_{\mathrm g}L}
\end{equation}
in place of $\exp\{CL\log(2+L)\}$. In particular, for $|j|\le L$,
\[
 \|\chi\circ T_R^j\|_{\mathrm{BV}}
 +\|h_R\circ T_R^j\|_{\mathrm{BV}}
 +\|\eta_R\circ T_R^j\|_{\mathrm{BV}}\le B_L,
 \qquad
 \|g\circ T_R^j\|_{\mathrm{BV}}\le C\|g\|_{C^1}B_L.
\]
The same exponential bound holds for every binary section-membership
word of length at most $2L+2$, extended by zero off its domain.
The constants are independent of the moving rectangle coefficients.
\end{theorem}
\begin{proof}
Retain the first-admissible-flight graph
\eqref{eq:polynomial-flight-graph}, including its no-earlier-hit
conditions, incidence signs and section-membership tests. Its degree
is bounded independently of $L$, whereas its variable and test counts
are $O(L)$. The number of label and membership alternatives is
exponential in $L$. Fixing one initial coordinate and a superlevel of
a bounded record component adds only a bounded number of variables
and tests. Lemma~\ref{lem:linear-sign-complexity} therefore bounds
the number of interval components of every projected slice superlevel
by $Ce^{CL}$. A finite angular atlas and its seam contributions change
only the constants.

For a bounded scalar function $u$ on an interval, with $|u|\le M$,
one-dimensional coarea gives
\[
 |Du|\le\int_{-M}^{M}\operatorname{Per}\{u>t\}\dd t.
\]
A superlevel with at most $b$ interval components has perimeter at most
$2b$, with the endpoint convention incorporated by zero extension.
Thus the slice variation is at most $4Mb$. Integrating in the other
initial coordinate proves both first distributional derivative bounds.
The statement for binary words follows in exactly the same way by
adjoining all their membership tests to the finite graph; it does not
require differentiating the word's boundary pointwise.

For smooth $g$, cover the angular circle by the same fixed regular
charts and apply the BV chain and product inequalities to the bounded
coordinate record $\chi\circ T_R^j$. The trace terms at the finitely
many chart seams have the same component bound. This gives the stated
$C^1$ composition estimate. Negative iterates use the unique regular
lifted fiber with the terminal state fixed, as in the original graph.
All coefficient dependence enters the sign conditions, to which
\eqref{eq:finite-binomial-exponential} applies uniformly.
\end{proof}

\subsection{The regularized tower with one logarithm}
Use $\mathsf L_{R,z,\xi}$ and $\mathcal E_{R,z,\xi}$ from
Section~\ref{sec:peripheral-return-phases}. Set
$M_f=\|f\|_\infty$ and $V_f=\|f^0\|_{\mathrm{BV}}$ for a section
function. The constants $a_{\mathrm t},c_{\mathrm t}>0$ below are chosen
so that the cumulative-return tail reads
\begin{equation}\label{eq:path-cumulative-tail-convention}
 \nu_R^*\{N_{m,R}>L\}\le C e^{a_{\mathrm t}m-c_{\mathrm t}L}.
\end{equation}
The one-return tail may have different positive constants.

\begin{proposition}[Exponential regularization budget]
\label{prop:exponential-tower-regularization}
For $L\ge1$ and $0<\epsilon<1/2$, the finite-age approximant of
Lemma~\ref{lem:exact-peripheral-lift}, denoted here by $Q$, satisfies
\begin{align}
 \|Q-\mathsf L_{R,z,\xi}f\|_1
   &\le C\{L\epsilon V_f+M_f e^{-cL}\},
       \label{eq:exponential-tower-error}\\
 \|Q\|_\infty&\le M_f,\qquad
 \|Q\|_{\mathrm{BV}}
   \le C M_f e^{CL}(1+\epsilon^{-1}+|z|).
       \label{eq:exponential-tower-BV}
\end{align}
Its squared $L^2$ error is bounded by $2M_f$ times the right side
of \eqref{eq:exponential-tower-error}. These assertions hold for every
real return phase $\xi$.
\end{proposition}
\begin{proof}
The age truncation, initial smoothing, exact top convention and
multiplication by $b_{R,\xi}$ are unchanged. In particular their
$L^1$ errors are precisely those of
Proposition~\ref{prop:tower-bv-approximation} and
Lemma~\ref{lem:exact-peripheral-lift}. For age $j<L$, the level
indicator is a binary backward membership word. Apply
Theorem~\ref{thm:exponential-finite-record} to that word, the smoothed
initial function evaluated at $T_R^{-j}$, and the accumulated phase.
The product and chain inequalities bound the variation by
$CM_f e^{Cj}(1+\epsilon^{-1}+j|z|)$. Summing the ages and absorbing
polynomial factors in $e^{CL}$ gives \eqref{eq:exponential-tower-BV}.
The age levels are disjoint, so the supremum bound is $M_f$, not $LM_f$.
Multiplication by $b_{R,\xi}$ costs only its uniform BV norm. Finally
both the lift and its approximant have modulus at most $M_f$, giving
the $L^2$ assertion. No regularity of the full lift is presumed.
\end{proof}

\begin{theorem}[A single-logarithm joint return estimate]
\label{thm:single-log-return-defect}
Put
\[
 \rho=(|z|^2+\xi^2)^{1/2},\qquad
 \Lambda=1+\log(H/\rho),\qquad H\ge2,
       \quad \xi\in[-\pi,\pi].
\]
There are uniform $a,b,\rho_0>0$ such that $0<\rho\le\rho_0$
and $\rho^2\Lambda\le a$ imply
\begin{align}
 |q|=1,\quad \|q^0\|_{\mathrm{BV}}\le H
 &\quad\Longrightarrow\quad
 \|\mathcal E_{R,z,\xi}q\|_{L^1(\nu_R^*)}\ge b\rho^2,
       \label{eq:single-log-circle}\\
 \|f\|_{L^2(\nu_R^*)}=1,\quad M_f+V_f\le H
 &\quad\Longrightarrow\quad
 \|\mathcal E_{R,z,\xi}f\|_{L^2(\nu_R^*)}
                    \ge b\rho^2/\Lambda.
       \label{eq:single-log-vector}
\end{align}
In the second assertion $f$ may vanish. The domain contains the pure
spectral direction and the shifted-frequency cancellation line.
\end{theorem}
\begin{proof}
Write $\sigma=c_*\xi$ and $w=z-\sigma e_3$. On a sufficiently small
fixed joint ball,
\[
 c_1\rho\le t=(|w|^2+\operatorname{dist}(\sigma,2\pi\Z)^2)^{1/2}
                         \le c_2\rho.
\]
This follows from the invertible linear frequency change in
\eqref{eq:peripheral-frequency-shift}; it does not divide by $|w|$.

For a unit section phase the exact lift has modulus one on the full
collision cylinder. Choose an $L^1$ error $\eta=\epsilon_0\rho^2$,
with $\epsilon_0>0$ fixed later. In
Proposition~\ref{prop:exponential-tower-regularization}, take
$L\le C(1+|\log\eta|)$ large enough for the tail and
$\epsilon=\eta/(CLH)$. Circle truncation, with the same coarea threshold
as in the proof of Theorem~\ref{thm:joint-return-spectral-defect},
gives a unit collision phase $v$ with
\[
 \|v-\mathsf L_{R,z,\xi}q\|_1\le C\eta,\qquad
 1+\log(2+\|v\|_{\mathrm{BV}})
               \le C(1+|\log\epsilon_0|)\Lambda.
\]
Indeed the logarithm in \eqref{eq:exponential-tower-BV} is at most
$CL+\log H+\log(1+\epsilon^{-1})+C$, and
$\log L\le C(1+|\log\eta|)$. Apply
Lemma~\ref{lem:joint-collision-phase-defect}, and use the exact
$L^1$ defect identity \eqref{eq:peripheral-lift-defect}. The change of
defect costs at most twice the approximation error. Choose
$\epsilon_0$ small enough to absorb it and then reduce the joint
smallness constants. This proves \eqref{eq:single-log-circle}.

For a normalized complex section function, set
$Q_*=\mathsf L_{R,z,\xi}f$ and $S=\|Q_*\|_2$. The exact norm
identity gives $\sqrt{c_*}\le S\le H$. If
$e=\|\mathcal E_{R,z,\xi}f\|_2$, the collision defect of $Q_*/S$
is at most $e$. Choose
\[
 \eta=\epsilon_0\rho^2/\Lambda,\qquad
 L\le C(1+\log(H/\eta)),\qquad
 \epsilon=\eta^2/(CH^2L),
\]
with the fixed constants large enough to make the normalized $L^2$
error at most $\eta$. Normalizing the approximant gives $\|v\|_2=1$
and
\begin{align*}
 \|v-Q_*/S\|_2&\le2\eta,\\
 1+\log(2+\|v\|_\infty+\|v\|_{\mathrm{BV}})+|\log t|
       &\le C(1+|\log\epsilon_0|)\Lambda.
\end{align*}
The last inequality follows from
$\log(H/\eta)=\log H+2\log(1/\rho)+\log\Lambda+
\log(1/\epsilon_0)$ and $\log\Lambda\le\Lambda$.
Proposition~\ref{prop:joint-collision-vector-defect} consequently gives
\[
 e+4\eta\ge\frac{b_0c_1^2\rho^2}
                         {C(1+|\log\epsilon_0|)\Lambda}.
\]
Since $\epsilon_0(1+|\log\epsilon_0|)\to0$, choose $\epsilon_0$
first so that the approximation term is absorbable. Only then reduce
$a,\rho_0$ to meet the collision budget. This proves
\eqref{eq:single-log-vector} with no circular accuracy choice.
\end{proof}

\begin{corollary}[Peripheral annuli and compressed resolvents]
\label{cor:single-log-compressed}
For $H\le H_0n^\zeta$, the preceding theorem applies on
\[
 2n^{-99/200}\le|z|\le n^{-2/5},\qquad
 |\xi|\le a_0/\sqrt{\log(2+n)}
\]
for sufficiently small fixed $a_0>0$ and sufficiently large $n$.
Its complex-function lower bound is
$c(|z|^2+\xi^2)/\log(2+n)$. The rescaled physical radii are still
$2n^{1/200}$ and $n^{1/10}$.

For the exact section-grid compression $A_{R,h,z}$ in
Section~\ref{sec:compressed-return-resolvents}, let
$H_h=C/h$ and $\Lambda_h=1+\log(H_h/\rho)$. On the joint region
$0<\rho\le\rho_0$, $\rho^2\Lambda_h\le a$,
\begin{equation}\label{eq:single-log-compressed-resolvent}
 \|(e^{i\xi}-A_{R,h,z})^{-1}\|
                          \le C\Lambda_h^2/\rho^4.
\end{equation}
The same bound, with twice the constant, holds when the resolvent
parameter lies within $c\rho^4/\Lambda_h^2$ of $e^{i\xi}$.
\end{corollary}
\begin{proof}
On the annulus $\Lambda\le C\log(2+n)$, and
$|z|^2\Lambda\to0$ uniformly. Choosing $a_0$ small controls the
spectral contribution to the budget. For a normalized grid vector,
Lemma~\ref{lem:section-grid-reconstruction} gives the budget $H_h$.
If its compressed residual is $\varepsilon\le1$,
Lemma~\ref{lem:compression-defect-comparison} bounds the squared
physical defect by $3\varepsilon$. Theorem~\ref{thm:single-log-return-defect}
therefore gives $\varepsilon\ge c\rho^4/\Lambda_h^2$. Residuals
larger than one satisfy the same lower bound after changing $c$.
This is a smallest-singular-value estimate, hence implies
\eqref{eq:single-log-compressed-resolvent}. Moving the scalar resolvent
parameter by at most half that lower bound preserves half the singular
value by the triangle inequality. No normality assumption is used.
\end{proof}

\begin{proposition}[Improved finite-time consistency]
\label{prop:exponential-compression-consistency}
The finite-time comparison of
Proposition~\ref{prop:compressed-finite-time-comparison} improves to
\begin{equation}\label{eq:exponential-compression-consistency}
 \left|\int_{Y_R^*}e^{-iz\cdot U_{n,R}}\dd\nu_R^*
                       -\langle1,A_{R,h,z}^{\,n}1\rangle\right|
 \le Cn\left\{e^{(a_{\mathrm t}n-c_{\mathrm t}L)/2}
          +\sqrt{h(1+|z|)e^{CL}}\right\},\qquad L\ge n.
\end{equation}
The Fourier sign and Hilbert pairing are those of that proposition;
$U_{n,R}=J_{n,R}-n\bar G_R$ is the actual centered record.
\end{proposition}
\begin{proof}
For each $j\le n$, restrict the actual $j$-return phase to
$\{N_{j,R}\le L\}$. Resolve this event by its binary section word
up to the last collision, and use
Theorem~\ref{thm:exponential-finite-record}. Products of word indicators
and the phase chain rule give variation at most
$C(1+|z|)e^{CL}$; the number of words and all factors polynomial in
$L$ are absorbed in the exponential. The discarded part has $L^2$
norm at most $Ce^{(a_{\mathrm t}n-c_{\mathrm t}L)/2}$.
Cellwise BV approximation on the fitted grid gives the second term
inside braces. Substitute these two estimates in the exact projection
power telescoping identity used in the cited proposition and sum
its $n$ terms. Neither the original law nor the matrix entries are
truncated in \eqref{eq:exponential-compression-consistency}.
\end{proof}

These estimates improve the reconstruction and consistency constants.
They do not assert a mesh-uniform power bound for the uncompressed
operator: the consistency scale and the local resolvent scale must
still be reconciled, and the remaining peripheral arcs must be treated.
First variation of a collision record is also distinct from the
second derivative sum of the raw densities in Theorem I.
